\section{Introduction} \label{sec:intro}

The Vera C.\ Rubin Observatory will conduct the 10-year Legacy Survey of Space and Time (LSST; \citealt{ivezic2019lsstfe7}), imaging roughly 18,000~deg$^{2}$ of the southern sky in six broadband, optical--NIR filters ($ugrizy$).
Achieving the LSST science goals requires calibrating the instrument with high precision and uniformity so that raw detector readouts can be transformed into science-ready images and catalogs.
Errors or non-uniformities in image calibration can introduce survey-scale photometric and astrometric systematics that propagate into downstream measurements and may be degenerate with astrophysical or cosmological signals \citep{esteves2023photometry-3c8}.
Key challenges to delivering this level of calibration include:
\begin{enumerate}
    \item \textbf{Instrument complexity:} The LSST Camera (LSSTCam) is the largest digital camera built for astronomy.
    Its 3.2-gigapixel focal plane consists of 205 fully depleted, backside-illuminated CCDs from two vendors (Teledyne e2v and International Technologies Laboratory, hereafter E2V and ITL).
    Detector-to-detector and amplifier-level variations require calibration and processing algorithms robust to heterogeneous sensor behavior.
    \item \textbf{Data volume and cadence:} LSSTCam typically acquires 30~s exposures with $\leq$3~s readout, yielding about 1000 exposures per night and nightly data rates of approximately 10~TB.
    \item \textbf{Calibration illumination and optics:} The calibration system (dome flats, LEDs, monochromator, and the Collimated Beam Projector or CBP) must provide uniform and repeatable illumination over a large focal plane through a fast ($f/1.2$) optical system.
    \item \textbf{Coupled instrumental effects:} Instrument artifacts must be corrected in the reverse chronological order from their origin in the photon transfer chain, yet the calibration data to model those artifacts are often obtained and combined in a forward (chronological) way.
    Many effects are partially degenerate, requiring iterative or ``bootstrapped'' approaches to create self-consistent calibration products.
    \item \textbf{Evolving conditions and cadence:} Changing environmental conditions, multi-epoch calibration datasets, and evolving observing strategies require calibration production to be robust to non-uniform operations.
    The calibrations for DP2 used datasets collected before the dome was completed, before installation of thermal regulation systems, and before finalization of the precise detector operating mode.
    \item \textbf{Verification, provenance, and operational reproducibility:} Calibration products must be accompanied by quantitative diagnostics for rapid data quality assessment, cross-time and cross-version comparisons, and clear acceptance and certification decisions for downstream processing.
    Products must be reproducible, versioned, and continuously validated as instruments and algorithms evolve.
\end{enumerate}

The LSST Data Management Science Pipelines \citep{2025rubn.rept...32N} implement this functionality using production pipelines and supporting tools to generate, verify, and publish calibration products for routine processing.
The design benefits from lessons learned in LSST Data Preview 1 (DP1) \citep{2025rubn.rept...31N} and precursor surveys, including the Hyper Suprime-Cam Subaru Strategic Program \citep{miyazaki2012hyper-103,aihara2022third-22a,bosch2017hyper-e24} and the Dark Energy Survey \citep{collaboration2005dark-0fe,diehl2012dark-a58,morganson2018dark-dc6}.

The first end-to-end demonstration of Rubin data processing on a substantial, internally consistent LSSTCam dataset is \emph{Data Preview 2} (DP2).
DP2 covers about $2{,}500~\mathrm{deg}^{2}$ in all six filters and includes commissioning and pre-survey operations data.
This paper describes the calibration-production algorithms and the operational approach used to generate DP2 calibration products.

Our goals are to (1) make the DP2 calibration products interpretable to science users, (2) quantify our confidence in the quality of released data products, (3) provide a reference for future cross-survey synergies and science analyses, (4) make the processing reproducible for pipeline developers and operators, and (5) provide a reference point for expected improvements on the path to the first LSST public data release.

This paper is organized as follows.
\S\ref{sec:data} introduces the observatory hardware and instrumentation that supplied the inputs for DP2 calibrations.
\S\ref{sec:processing} presents the physical detector model and the calibration and instrument signature removal (ISR) pipeline, including algorithms and data products used in production.
\S\ref{sec:discussion} discusses calibration-epoch differences and residual systematics relevant to DP2.
\S\ref{sec:verify-accept-certify} describes how calibration products are verified, accepted, certified, and deployed.
\S\ref{sec:conclusion} summarizes the DP2 calibration quality and outlines expected improvements toward Data Release~1 (DR1).
